\documentclass[10pt,a4paper]{amsart}
\usepackage[margin=19mm]{geometry}
\usepackage{amsmath,amssymb,mathtools}
\usepackage[scr=rsfs,cal=euler]{mathalfa}
\usepackage{xfrac}
\usepackage[unicode,hidelinks,bookmarks=false]{hyperref}
\newcommand{\ie}{i.e.\ }
\newcommand{\cf}{cf.\ }
\newcommand{\resp}{respectively\ }
\newcommand{\Subsection}{\subsection}
\providecommand{\nobreakdash}{\nobreak\hbox{-}}
\setlength{\emergencystretch}{2em}
\multlinegap0pt
\usepackage{amssymb}
\usepackage{dsfont}
\usepackage[shortlabels]{enumitem}
\usepackage[normalem]{ulem}
\newtheorem{lemma}{Lemma}
\DeclareMathOperator{\Pol}{Pol}
\title{The ordered F-system for the F-method and differential symmetry breaking operators for $(GL(3,\mathbb{R}), GL(2,\mathbb{R}))$}
\author{Jonathan Ditlevsen, Toshihisa Kubo, V{\' i}ctor P{\'e}rez-Vald{\'e}s}
\date{Source: arXiv:2607.27956v1 (CC BY 4.0)}
\begin{document}
\maketitle

\noindent\emph{Source and licence.} The ordered F-system for the F-method and differential symmetry breaking operators for $(GL(3,\mathbb{R}), GL(2,\mathbb{R}))$. Version arXiv:2607.27956v1 (CC BY 4.0), distributed by arXiv under the Creative Commons Attribution 4.0 International licence, http://creativecommons.org/licenses/by/4.0/, as recorded in the arXiv licence metadata for this exact version and shown on the abstract page https://arxiv.org/abs/2607.27956v1. Full licence notice: \texttt{licenses/arxiv-2607.27956-CC-BY-4.0.txt}.\par\smallskip

\noindent\textit{Editorial scope.} Counted unit: the article's lemma lem:eDD (source0.tex lines 8548--8553). The article's complete original proof of it is delivered with it (source0.tex lines 8555--8641, one \texttt{proof} environment of 87 lines). No mathematics is written, repaired or re-derived: the statement and the proof are the article's own lines, in the article's own notation, and no retained line is substituted or re-flowed.  Retained uncounted as the source-local context the proof needs: lines 8404--8408.  Nothing was omitted from any retained span, so the package carries no omission record.  No retained line contains a citation, so the wrapper carries no bibliography and no bibliography entry.  No retained line contains a figure environment, an image-inclusion command or drawing code, so no image is delivered and the package declares no asset dependency.  Editorial formatting only, in addition: an AMS \texttt{amsart} preamble that restates the article's own declarations and macros used by the retained lines, namely the environments \texttt{lemma} on separate counters, together with the standard packages those lines need.  The retained environments print in this document as Lemma 1 in order of appearance, whereas the article prints the same items with the numbering of its own full document.  The complete native source of the article as distributed in the arXiv e-print for this exact version is bundled unchanged as \texttt{sources/206377/source0.tex}.
\begin{equation}\label{eqn:eJk}
e^{\tilde{J}(\alpha_1, \alpha_2)}t^k
= \sum_{n=0}^{\infty}\frac{(k-\alpha_1)_n(k-\alpha_2)_n}{2^nn!}t^{n+k}
= \sum_{n=0}^{\min(\alpha)-k}\frac{(k-\alpha_1)_n(k-\alpha_2)_n}{2^nn!}t^{n+k}.
\end{equation}

\begin{lemma}\label{lem:eDD}
As elements of $GL(\Pol_{\min(\alpha)}[t])$ we have the following identity:
\begin{equation*}
e^{\tilde{J}(\alpha_1,\alpha_2+1)}e^{-\tilde{J}(\alpha_1,\alpha_2)}=1+\tfrac{t}{2}(\alpha_1-\vartheta_t).
\end{equation*}
\end{lemma}

\begin{proof}
Take $t^k \in \Pol_{\min(\alpha)}[t]$. We wish to prove
\begin{equation*}
e^{\tilde{J}(\alpha_1,\alpha_2+1)}e^{-\tilde{J}(\alpha_1,\alpha_2)}t^k
=\big(1+\tfrac{t}{2}(\alpha_1-\vartheta_t)\big)t^k.
\end{equation*}
Clearly, we have 
\begin{equation*}
\big(1+\tfrac{t}{2}(\alpha_1-\vartheta_t)\big)t^k 
= \big(1+\tfrac{t}{2}(\alpha_1- k)\big)t^k.
\end{equation*}
Thus we aim to show
\begin{equation*}
e^{\tilde{J}(\alpha_1,\alpha_2+1)}e^{-\tilde{J}(\alpha_1,\alpha_2)}t^k
= \big(1+\tfrac{t}{2}(\alpha_1- k)\big)t^k.
\end{equation*}

In the following, we simply write finite sums as infinite sums as in \eqref{eqn:eJk}.
Then, by \eqref{eqn:eJk}, we have 
\begin{equation*}
e^{-\tilde{J}(\alpha_1, \alpha_2)}t^k
= \sum_{n=0}^{\infty}(-1)^n\frac{(k-\alpha_1)_n(k-\alpha_2)_n}{2^nn!}t^{n+k}.
\end{equation*}
Thus, $e^{\tilde{J}(\alpha_1,\alpha_2+1)}e^{-\tilde{J}(\alpha_1,\alpha_2)}t^k$ 
is computed as
\begin{multline}
e^{\tilde{J}(\alpha_1,\alpha_2+1)}e^{-\tilde{J}(\alpha_1,\alpha_2)}t^k
=\\
\sum_{m=0}^{\infty}\sum_{n=0}^{\infty}
(-1)^n\frac{(k-\alpha_1)_n(k-\alpha_2)_n(n+k-\alpha_1)_m(n+k-\alpha_2-1)_m}{2^{n+m}n!m!} t^{n+m+k}. \label{eqn:eJeJ}
\end{multline}
By the property $(x)_\ell(x+r)_k = (x)_{r+k}$,
the expression \eqref{eqn:eJeJ} can be reduced to
\begin{align}
\eqref{eqn:eJeJ}
&=\sum_{m=0}^{\infty}\sum_{n=0}^{\infty}
(-1)^n\frac{(k-\alpha_1)_n(k-\alpha_2)_n(n+k-\alpha_1)_m(n+k-\alpha_2-1)_m}{2^{n+m}n!m!} t^{n+m+k} \nonumber \\[3pt]
&=t^k \sum_{\ell=0}^\infty \frac{(k-\alpha_1)_\ell}{2^\ell}t^\ell
\left(\sum_{n+m=\ell} (-1)^n\frac{(k-\alpha_2)_n(n+k-\alpha_2-1)_m}{n!m!}\right)\nonumber\\[3pt]
&=\frac{t^k}{k-\alpha_2-1}
 \sum_{\ell=0}^\infty \frac{(k-\alpha_1)_\ell(k-\alpha_2-1)_\ell}{2^\ell}t^\ell
\left(\sum_{n+m=\ell} \frac{(-1)^n}{n!m!}(n+k-\alpha_2-1)\right)\nonumber\\[3pt]
&=\frac{t^k}{k-\alpha_2-1}
 \sum_{\ell=0}^\infty \frac{(k-\alpha_1)_\ell(k-\alpha_2-1)_\ell}{2^\ell \ell!}t^\ell
\left(\sum_{n=0}^\ell(-1)^n\binom{\ell}{n}(n+k-\alpha_2-1)\right).\label{eqn:eJeJ2}
\end{align}
For the inner sum $\sum_{n=0}^\ell(-1)^n\binom{\ell}{n}(n+k-\alpha_2-1)$ 
in \eqref{eqn:eJeJ2}, observe that 
\begin{alignat*}{2}
\sum_{n=0}^\ell(-1)^n\binom{\ell}{n}
&= (1-x)^\ell\big\vert_{x=1} 
&&= 
\begin{cases}
1 & \text{if $\ell = 0$},\\
0 & \text{otherwise},
\end{cases}\\[5pt]
\sum_{n=1}^\ell(-1)^n\binom{\ell}{n}n
&= \frac{d}{dx}\bigg\vert_{x=1}(1-x)^\ell 
&&= 
\begin{cases}
-1 & \text{if $\ell = 1$},\\
0 & \text{otherwise}.
\end{cases}\\[3pt]
\end{alignat*}
Thus, the sum $\sum_{n=0}^\ell(-1)^n\binom{\ell}{n}(n+k-\alpha_2-1)$ is 
evaluated to
\begin{align*}
\sum_{n=0}^\ell(-1)^n\binom{\ell}{n}(n+k-\alpha_2-1)
=
\begin{cases}
(k-\alpha_2-1) & \text{if $\ell=0$},\\[5pt]
-1 & \text{if $\ell=1$},\\[3pt]
0 & \text{otherwise}.
\end{cases}
\end{align*}
Hence, we have 
\begin{align*}
\eqref{eqn:eJeJ2}
&=\frac{t^k}{k-\alpha_2-1}
\sum_{\ell=0}^\infty \frac{(k-\alpha_1)_\ell(k-\alpha_2-1)_\ell}{2^\ell \ell!}t^\ell
\left(\sum_{n=0}^\ell(-1)^n\binom{\ell}{n}(n+k-\alpha_2-1)\right)\\[5pt]
&=\frac{t^k}{k-\alpha_2-1}
\big( (k-\alpha_2-1)-\tfrac{t}{2} (k-\alpha_2-1)(k-\alpha_1)\big)\\[5pt]
&=\big(1+\tfrac{t}{2}(\alpha_1-k)\big)t^k.
\end{align*}
This completes the proof.
\end{proof}

\end{document}
